\section{Physical boundary layers on the image collision space}
\label{sec:v49-physical-layers}
We estimate the true incidence and clearance defects without continuing
an orbit across a physical singularity. The clearance of flight $j$ is
read at collision $j+1$, from the preceding flight. Its thin positive
upper envelope then has boundaries transverse to the stable cone.
This gives a local bound for the actual defective histories, not for
histories with an inserted, deleted or relabelled collision.

The table in this section is the original circular family. Write
$r=R\alpha$ for arclength, $p=\sin\varphi$, and
$c(x)=\cos\varphi$. The section mass remains
$c=91/(10000\pi)$ when no argument is displayed. To avoid confusion,
the incidence below is denoted by $c_j$, as in
Section~\ref{sec:v45-critical-collars}. Choose
\[
 s_0=1/100,\qquad
 \ell_* =\sqrt{(53/100)^2-(48/100)^2}>0.
\]
Enlarging a constant permits any fixed multiple of a width smaller
than $s_0$. All multipliers use the density convention and local
common spaces of Section~\ref{sec:action-norm-details}.

\subsection{A one-flight geometric upper envelope}
At an image collision with its disk centered at zero, put
\[
 q_R=Rn_\alpha,\quad
 v^- =\cos\varphi\,n_\alpha-\sin\varphi\,t_\alpha,\quad
 (v^-)^\perp=\sin\varphi\,n_\alpha+\cos\varphi\,t_\alpha.
\]
Thus $v^-$ points backwards along the incoming flight. Let $\mathcal D$
be the finite set of nonzero lattice centers of norm at most
$\tau_{\max}+2R_++s_0$. For $d\in\mathcal D$ define
\begin{equation}\label{eq:v49-backward-distance}
 L_d=(d-q_R)\cdot v^-,\qquad
 Z_d=(d-q_R)\cdot(v^-)^\perp.
\end{equation}
Here $L_d$ is a geometric distance, not a reconstruction band or a
collision count. Define the two indicators
\begin{align}
 g_{R,s}(x)&=\one_{\{\cos\varphi<s\}},\notag\\
 b_{R,s}^-(x)&=\one_{\bigcup_{d\in\mathcal D}
           \{L_d>0,\ R\le |Z_d|\le R+s\}}(x).
                                      \label{eq:v49-layer-indicators}
\end{align}
The second set is an upper envelope. It does not require that the
perpendicular foot occur before the preceding collision, and it is
not asserted to be the actual clearance set.

\begin{lemma}[Transverse envelopes for actual physical defects]
\label{lem:v49-physical-envelope}
For every regular flight with clearance $\Delta_j<s<s_0$,
\begin{equation}\label{eq:v49-clearance-inclusion}
 b_{R,s}^-(T_R^{j+1}x)=1.
\end{equation}
On a fixed angular chart the boundary curves of the clearance
envelope are uniformly $C^2$, and are transverse to stable graphs.
Every homogeneous stable curve intersects the envelope in a uniformly
bounded number of intervals, each of length at most $Cs$. The same
conclusion holds for $g_{R,s}$, including the homogeneity strips
accumulating at grazing. Partition and matching constants are
independent of $s$ and $R$.
\end{lemma}
\begin{proof}
The endpoints of a physical flight are on its two incident disks.
Their distances from any other disk are at least $3/50$. Since
$s<s_0$, a nonincident disk at clearance less than $s$ has its nearest
point strictly inside the segment. Reading this segment backwards
from its terminal collision, the foot therefore has $L_d>0$, the
supporting line has distance $|Z_d|\in[R,R+s]$, and $d$ belongs to
$\mathcal D$. This proves \eqref{eq:v49-clearance-inclusion}. The
actual previous disk can be left in the upper-envelope candidate
list: an extra upper-comparison set causes no identification of labels.

On the indicated strips, $|d-q_R|\ge1-R\ge53/100$ and
$|Z_d|\le R+s_0\le48/100$, so $L_d\ge\ell_*$. The sign condition
$L_d>0$ creates no additional boundary meeting these strips.
Direct differentiation in $(r,\varphi)$ gives
\begin{equation}\label{eq:v49-transverse-derivatives}
 \partial_r Z_d=-\cos\varphi-L_d/R,
 \qquad \partial_\varphi Z_d=L_d.
\end{equation}
Consequently a level $Z_d=z$ in a strip is a graph with slope
\[
 \frac{d\varphi}{dr}=R^{-1}+\frac{\cos\varphi}{L_d}>0.
\]
Its slope and first two derivatives have common upper bounds, since
$L_d\ge\ell_*$ and $\mathcal D$ is finite. A stable graph has negative
slope bounded away from zero and infinity. Along it,
\[
 \frac{dZ_d}{dr}=-\cos\varphi-L_d/R+L_d\varphi'(r)
                  \le-\ell_*/R_+.
\]
Thus a change of $Z_d$ by $s$ occupies length at most $Cs$.
The four levels $\pm R,\pm(R+s)$ and fixed angular chart cuts form
an admissible partition. Their number of components in each
homogeneity strip is uniformly bounded: at a horizontal strip edge
the level equation is a trigonometric equation of bounded degree in
$\alpha$, with a fixed number of candidates. Refining by these finitely
many arcs gives simply connected cells. Uniform transversality gives
length $O(\delta)$ in a $\delta$-neighborhood of every cell boundary.
Neither bound depends on the separation of the parallel level curves.

For the incidence envelope, $\cos\varphi<s$ is the union of two
horizontal bands of angular width $\arcsin s\le2s$. Stable slopes
are bounded away from zero, so each intersection has length $Cs$.
A homogeneous stable curve is already inside one homogeneity strip.
Most such strips are either entirely included or entirely excluded;
only the two horizontal thresholds require new cuts. The number of
new cuts per strip is bounded, even when the threshold approaches
$\pm\pi/2$. No inverse change from $p$ to $\varphi$ is used at grazing.
Matching intersections of two stable graphs moves a transverse cut
by at most a constant times their graph distance. If the thin band is
smaller than that distance, its whole intersection has that same
unmatched-length bound. This proves the stated uniformity.
\end{proof}

\subsection{Strong multiplication and a weak gain}
\begin{proposition}[Physical thin-layer multipliers]
\label{prop:v49-physical-multiplier}
For $a_{R,s}=g_{R,s}$ or $b_{R,s}^-$, uniformly on each local common
collision space,
\begin{equation}\label{eq:v49-physical-multiplier}
 \|M_{a_{R,s}}h\|_{\mathcal B}\le C\|h\|_{\mathcal B},\qquad
 |M_{a_{R,s}}h|_w\le Cs^{1/8}\|h\|_{\mathcal B}.
\end{equation}
These assertions hold for the physical multiplication operation on
the strong completion. They do not assert a small strong norm.
\end{proposition}
\begin{proof}
Use the partition in Lemma~\ref{lem:v49-physical-envelope} and the
piecewise multiplier criterion \cite[Lemma 3.3(b), p.~13]{DPZ}, with
the pinned numbering specified in the bibliography. On each partition
cell the indicator is constant. The cell-count, intersection and
boundary-neighborhood hypotheses have just been verified uniformly.
The fixed exponent choices of Section~\ref{sec:action-norm-details}
satisfy the multiplier restrictions.

For the weak estimate, restrict a weak stable test to its at most
$N$ retained intervals $U$. Each has length at most $Cs$. The strong
stable length normalization, with $\varsigma=1/8$, gives
\[
 \left|\int_U h\psi\,dm_U\right|
              \le C\|h\|_s|U|^{\varsigma}.
\]
Summation proves the second estimate. In the strong unstable
comparison the unmatched lengths are $O(\delta)$, including when
$s<\delta$, and cost $C\delta^{\varsigma}$. This is admissible
because $\gamma<\varsigma$. On retained pieces the multipliers
agree. No inverse-width derivative is paid. The multiplier theorem
extends this physical product from smooth densities to the strong
completion. Equivalently, smooth averages of the transverse level
indicators have uniform strong bounds and converge in the
strong-to-weak norm; their physical integrals converge by domination.
The local changes of arclength have uniformly bounded scale factors.
\end{proof}

\begin{theorem}[A marked physical-defect local upper bound]
\label{thm:v49-physical-marked-local}
Put $\alpha_*=1/16$. There is $C$ such that, for each fixed auxiliary
roof band $B_*\ge B_0$, there are $C_{B_*}<\infty$ and
$\rho_{B_*}<1$ for which
\begin{align}
 &m^2\nu\{K^{\rm c}_{m,R}=k,A_{m,R}=l,
       S_{m,R}\in J,\ a_{R,s}(T_R^j x)=1\}\notag\\
 &\qquad\le Cs^{\alpha_*}(|J|+B_*^{-1})
       +C_{B_*}(1+B_*|J|)m^2\rho_{B_*}^{m/2}.
                                      \label{eq:v49-physical-local}
\end{align}
This holds for either multiplier in
\eqref{eq:v49-layer-indicators}, every interval $J$, every exact
discrete target, $m\ge2$, $0\le j\le m$, $0<s<s_0$, and $R\in[0.45,0.47]$.
The leading constant is independent of $B_*$. The mark and width may
vary with $m$. A defect of flight $j$ is estimated with the clearance
mark at $j+1$, not at $j$.
\end{theorem}
\begin{proof}
The exact full-occupation pairing is
\[
 \ell(\mathscr T_{R,\theta}^{m-j}M_{a_{R,s}}
                    \mathscr T_{R,\theta}^{j}\nu).
\]
It has the physical meaning in
\eqref{eq:v48-marked-pairing}, now with the physical-layer multiplier.
On every fixed moving-peak chart,
Lemma~\ref{lem:v48-weak-interpolation} and
\eqref{eq:v49-physical-multiplier} give projection and terminal
amplitudes at most $Cs^{1/16}$. At least one chronological block has
length $m/2$. Expanding that block, and using the strong power bound
for the other one, gives
$Cs^{1/16}|\lambda|^{m/2}+C\rho^{m/2}$.
The four-frequency Gaussian peak bound integrates this to
$Cs^{1/16}m^{-2}+C\rho^{m/2}$.
The complement on a fixed outer band has an exponentially decaying
long block. All occupation resonances are retained. The positive
interval majorant has mass $|J|+2\pi/B_*$ and Fourier support in
$[-B_*,B_*]$. Its inversion gives
\eqref{eq:v49-physical-local}, exactly as in
Theorem~\ref{thm:v48-marked-strip-local}. Constants on the complement
may depend on $B_*$; none depends on $j$ or $s$.
\end{proof}

\paragraph{Physical interpretation.}
This is an estimate on the union of actual singular regimes, not a
transport across a grazing or competing-hit seam. Each trajectory
keeps its own $m$ collisions and all its original occupation decisions.
The source at a fixed return index is bounded above by the ambient
source $\nu/c$ with the same $(k,n,m)$ coefficient. No factor of $c$
is absorbed into a changed probability and no extra terminal visit is
added to $A_{m,R}=\sum_{j=0}^{m-1}\eta_R\circ T_R^j$.
